\documentclass{article}
\usepackage[utf8]{inputenc}
\usepackage{graphicx}
\usepackage{amsmath}
\usepackage{amssymb}
\usepackage{url}
\usepackage{tabularx}
\usepackage[colorlinks=true,allcolors=blue]{hyperref}
\usepackage{caption}
\usepackage{pdfpages}
\usepackage{multicol}
\usepackage[ngerman]{babel}
\usepackage{multicol}
\usepackage{geometry}	% [textwidth=17cm,textheight = 24.5cm]
% Um die Dokumentränder zu sehen: [showframe] vor {geometry} einfügen
\usepackage{bbold}
\usepackage{makecell}
\usepackage{float}
\usepackage{braket}
\usepackage{marvosym}
\usepackage{cancel}
\usepackage{xltabular}

%\usepackage{pascaltriangle}

\renewcommand{\cellalign}{tl}

\setlength{\columnsep}{1cm}
\setlength\parindent{0pt}

\newcommand{\diff}{\mathrm{d}}
\newcommand{\konst}{\mathrm{konst.}}
\newcommand{\lr}[1]{\left(#1\right)}
\newcommand{\lreck}[1]{\left[#1\right]}
\newcommand{\lrschweif}[1]{\left\{#1\right\}}
\newcommand{\dkon}[3]{\left(\frac{\partial #1}{\partial #2}\right)_{#3}}

\definecolor{gruen}{HTML}{038133}

\newcommand{\commenttext}[1]{\textcolor{gruen}{\text{#1}}}
\newcommand{\comment}[1]{{\color{gruen}{#1}}}


\begin{document}


\begin{center}
	\Large{\textbf{Analysis III Zusammenfassung}}\\
	\small{613348}\\
	\today
\end{center}

\tableofcontents


\newpage


\section*{Übersicht}
Das Modul Analysis III behandelt:
	\begin{itemize}
		\item \textbf{Differentialgleichungen}
		
		
		\item \textbf{Operatoren}
		
		
		\item \textbf{Fourier-Reihe/-Transformation}
		
		
		\item \textbf{Lebesgue-Integration}
	\end{itemize}


\newpage


\section{Fourier-Reihe/-Transformation}
\subsection{Fourier-Reihe}
\subsubsection{Trigonometrische Darstellung}
Die Fourier-Reihe ist die Darstellung einer periodischen Funktion $f(x)$ mit der Periodendauer $T$ in der \underline{diskreten} Basis $\{1, \cos \lr{ \frac{2\pi \cdot 1}{T} \cdot x }, \cos \lr{ \frac{2\pi \cdot 2}{T} \cdot x}, ..., \sin \lr{ \frac{2\pi \cdot 1}{T} \cdot x }, \sin \lr{ \frac{2\pi \cdot 2}{T} \cdot x }, ... \}$:
	\begin{equation}
		f(x) = \frac{a_0}{2} + \sum_{n=1}^\infty \lreck{ a_n \cdot \cos \lr{ \frac{2\pi \cdot n}{T} \cdot x } + b_n \cdot \sin \lr{ \frac{2\pi \cdot n}{T} \cdot x } }
	\end{equation}

mit den \textbf{Fourier-Koeffizienten}
	\begin{align}
		a_n &= \frac{2}{T} \int_T \diff x \, f(x) \cdot \cos \lr{ \frac{2\pi \cdot n}{T} \cdot x } \\
		b_n &= \frac{2}{T} \int_T \diff x \, f(x) \cdot \sin \lr{ \frac{2\pi \cdot n}{T} \cdot x },
	\end{align}

wobei $n \in \{0, 1, ...\}$.\\
Die Integration über die Periodendauer gilt dabei für jedes Intervall der Länge $T$, z.B. $[0,T]$, $[a,a+T]$, $[-T/2,T/2]$, etc.

\vspace{0.5cm}

{\color{blue}
\hrule
	\fbox{Beispiel}\\
	ges: Fourier-Reihe für die Rechteckfunktion\\
	\\
	Die Rechteckfunktion auf dem Intervall $[0,T]$ wird beschrieben durch die Funktion
		\begin{align*}
			f(x) =
			\begin{cases}
				1, & 1 \leq x \leq T/2 \\
				0, & T/2 < x \leq T
			\end{cases}.
		\end{align*}

	Die Fourier-Koeffizienten sind
		\begin{align*}
			\bullet \ a_0 &= \frac{2}{T} \int_0^{T/2} \diff x \, 1 \cdot 1 = \frac{2}{T} \cdot \frac{T}{2} = 1 \\
			\bullet \ a_n &= \frac{2}{T} \int_0^{T/2} \diff x \, 1 \cdot \cos \lr{ \frac{2\pi \cdot n}{T} \cdot x } = \frac{\cancel{2}}{\cancel{T}} \cdot \frac{\cancel{T}}{\cancel{2}\pi \cdot n} \int_0^{u(T/2) = \pi \cdot n} \diff u \, \cos(u) = \\
			&= \frac{1}{\pi \cdot n} \lreck{ \sin \lr{ \pi \cdot n } - \sin(0) } = 0 \\
			\bullet \ b_n &= \frac{2}{T} \int_0^{T/2} \diff x \, 1 \cdot \sin \lr{ \frac{2\pi \cdot n}{T} \cdot x } = \frac{\cancel{2}}{\cancel{T}} \cdot \frac{\cancel{T}}{\cancel{2}\pi \cdot n} \int_0^{u(T/2) = \pi \cdot n} \diff u \, \sin(u) = \\
			&= -\frac{1}{\pi \cdot n} \lreck{ \cos (\pi \cdot n) - \cos(0) } = \frac{ 1 - \cos(\pi \cdot n) }{ \pi \cdot n } =
			\begin{cases}
				\frac{2}{\pi \cdot n}, & \text{$n$ ungerade} \\
				0, & \text{$n$ gerade}
			\end{cases}
		\end{align*}
	
	Somit ist die Fourier-Reihe:
		\begin{align*}
			f(x) = \frac{1}{2} + \sum_{n=1}^\infty \frac{2}{\pi \cdot (2n-1)} \cdot \sin \lr{ \frac{2\pi \cdot (2n-1) }{T} \cdot x }
		\end{align*}
\hrule
}


\newpage


{\color{violet}
\hrule
	\fbox{Herleitung}\\
	geg: Periodische Funktion $f(x)$ mit der Periodendauer $T$\\
	zz: $f(x) = \frac{a_0}{2} + \sum_{n=1}^\infty \lreck{ a_n \cdot \cos\lr{ \frac{2\pi \cdot n}{T} \cdot x } + b_n \cdot \sin\lr{ \frac{2\pi \cdot n}{T} \cdot x } }$ mit den Fourier-Koeffizienten $a_n = \frac{2}{T} \int_T f(x) \cdot \cos \lr{ \frac{2\pi \cdot n}{T} \cdot x } \, \diff x$ und $b_n = \frac{2}{T} \int_T f(x) \cdot \sin \lr{ \frac{2\pi \cdot n}{T} \cdot x } \, \diff x$\\
	\\
	Strategie:\\
	1) Zeige, dass die Basis orthonormal und vollständig ist\\
	2) Entwickle die Funktion in der Basis\\
	3) Bestimme die Entwicklungskoeffizienten durch Integration\\
	\\
	Zeige die Orthonormalität, indem alle möglichen Skalarprodukte der Basis ausgerechnet werden. Die Basisvektoren lassen sich allgemein formulieren als
		\begin{align*}
			\sin\lr{ \frac{2\pi \cdot n}{T} x }, \ \cos\lr{ \frac{2\pi \cdot m}{T} x },
		\end{align*}
	
	mit $n \in \{1, 2, ... \}$ und $m \in \{0, 1, ...\}$.\\
	\\
	Wähle als Integrationsinterval oBdA $[-T/2,T/2]$.\\
	Somit sind die möglichen Skalarprodukte:
		\begin{align*}
			\bullet \ &\phantom{=} \int_{-T/2}^{T/2} \diff x \, \sin\lr{ \frac{2\pi \cdot n}{T} x } \cdot \cos\lr{ \frac{2\pi \cdot m}{T} x }
			\intertext{Verwende den trigonometrischen Zusammenhang $\sin(A) \cdot \cos(B) = \frac{1}{2} \lr{ \sin(A+B) - \sin(A-B) }$:}
			&= \frac{1}{2} \int_{-T/2}^{T/2} \diff x \, \sin\lr{ \frac{2\pi (n+m)}{T} x } - \frac{1}{2} \int_{-T/2}^{T/2} \diff x \, \sin\lr{ \frac{2\pi (n-m)}{T} x } = \\
			\intertext{$u$-Substitution des jeweiligen Sinusarguments liefert:}
			&= \frac{1}{2} \frac{T}{2\pi(n+m)} \int_{-u(T/2)}^{u(T/2)} \diff u \, \sin(u) - \frac{1}{2} \frac{T}{2\pi(n-m)} \int_{-u(T/2)}^{u(T/2)} \diff u \, \sin(u) = \\
			&= \frac{1}{2} \frac{T}{2\pi(n+m)} \cdot \lreck{ - \cos \lr{ \frac{2\pi (n+m)}{T} \frac{T}{2} } + \cos \lr{ \frac{-2\pi (n+m)}{T} \frac{T}{2} } } - \\
			&= - \frac{1}{2} \frac{T}{2\pi(n-m)} \cdot \lreck{ - \cos \lr{ \frac{2\pi (n-m)}{T} \frac{T}{2} } + \cos \lr{ \frac{-2\pi (n-m)}{T} \frac{T}{2} } } =
			\intertext{Aus der Symmetrie des Cosinus $\cos(x) = \cos(-x)$ folgt, dass sich die jeweils eingeklammerten Terme wegheben:}
			&= 0 \\
			\bullet \ &\phantom{=} \int_{-T/2}^{T/2} \diff x \, \sin\lr{ \frac{2\pi \cdot n}{T} x } \cdot \sin\lr{ \frac{2\pi \cdot m}{T} x } =...=
			\begin{cases}
				0, & n\neq m \\
				T/2, & n=m
			\end{cases}
		\end{align*}
	
	Analog für das Cosinus-Cosinus-Integral.\\
	\textcolor{red}{ERGÄNZEN Vollständigkeitsbeweis (vielleicht aber auch einfach hinnehmen)}


\newpage


	Die Entwicklung der Funktion $f(x)$ in der Basis hat die Form
		\begin{align*}
			f(x) &= \sum_{n=0}^\infty a_n \cdot \cos \lr{ \frac{2\pi n}{T} x } + \sum_{n=1}^\infty b_n \cdot \sin \lr{ \frac{2\pi n}{T} x } = \\
			&= a_0 + \sum_{n=1}^\infty \lreck{ a_n \cdot \cos\lr{ \frac{2\pi \cdot n}{T} \cdot x } + b_n \cdot \sin\lr{ \frac{2\pi \cdot n}{T} \cdot x } }.
		\end{align*}

	Die einzelnen Fourier-Koeffizienten folgen nun durch Multiplikation mit einer Basisfunktion, gefolgt von der Integration:
		\begin{align*}
			\bullet \ &\phantom{=} \int_T \diff x \, f(x) \cdot \sin \lr{ \frac{2\pi n}{T} x } = \\
			&= \int_T \diff \, x \lreck{ a_0 + \sum_{m=1}^\infty \lreck{ a_m \cdot \cos\lr{ \frac{2\pi \cdot m}{T} \cdot x } + b_m \cdot \sin\lr{ \frac{2\pi \cdot m}{T} \cdot x } } } \cdot \sin \lr{ \frac{2\pi \cdot n}{T} x } =
			\intertext{Aus der Orthonormalität der Basis folgt, dass lediglich der Term mit $b_m$ nicht verschwindet:}
			&= \sum_{n=1}^\infty b_n \cdot \delta_{nm} \cdot \frac{T}{2} = b_n \cdot \frac{T}{2} \\
			\Leftrightarrow b_n &= \frac{2}{T} \int_T \diff x \, f(x) \cdot \sin\lr{ \frac{2\pi n}{T}x } \qquad \checkmark
		\end{align*}
	
	\textcolor{red}{ERGÄNZEN cos}\\
	\\
	Die einfache Integration liefert
		\begin{align*}
			\bullet \, &\phantom{=} \int_T \diff x \, f(x) = a_0 \cdot T \\
			\Leftrightarrow a_0 &= \frac{1}{T} \int_T \diff x \, f(x).
		\end{align*}
	
	Damit $a_0$ dieselbe Definition wie $a_n$ mit $n\neq 0$ verwenden kann, wird lediglich der Faktor $\frac{1}{2}$ in der Reihenentwicklung ergänzt. $\qquad \blacksquare$
\hrule
}


\newpage


\subsubsection{Exponentialform-Darstellung}
Eine alternative Darstellung der Fourier-Reihe nutzt die komplexe Exponentialfunktion
	\begin{equation}
		f(x) = \sum_{q=-\infty}^{+\infty} c_q \cdot e^{i\frac{2\pi \cdot q}{T} \cdot x}
	\end{equation}

mit dem Entwicklungskoeffizienten
	\begin{equation}
		c_q = \frac{1}{T} \int_T f(x) \cdot e^{-i \frac{2\pi \cdot q}{T} x}.
	\end{equation}

\vspace{0.5cm}

{\color{violet}
\hrule
	\fbox{Herleitung}\\
	Strategie: Ausgehend von der trigonometrischen Darstellung, können Sinus und Cosinus durch die komplexe Exponentialfunktion ausgedrückt werden.\\
	\\
	Setze die Exponentialausdrücke für Cosinus und Sinus
		\begin{align*}
			\cos \lr{ \frac{2\pi \cdot n}{T} \cdot x } &= \frac{1}{2} \lr{ e^{i \frac{2\pi \cdot n}{T} x} + e^{-i \frac{2\pi \cdot n}{T} x} } \\
			\sin \lr{ \frac{2\pi \cdot n}{T} \cdot x } &= \frac{1}{2i} \lr{ e^{i \frac{2\pi \cdot n}{T} x} - e^{-i \frac{2\pi \cdot n}{T} x} }
		\end{align*}

	anstelle der Basisfunktionen in der Fourier-Reihe ein:
		\begin{align*}
			f(x) &= \frac{a_0}{2} + \sum_{n=1}^\infty \lreck{ a_n \cdot \frac{1}{2} \lr{ e^{i \frac{2\pi \cdot n}{T} x} + e^{-i \frac{2\pi \cdot n}{T} x} } + b_n \cdot \frac{1}{2i} \lr{ e^{i \frac{2\pi \cdot n}{T} x} - e^{-i \frac{2\pi \cdot n}{T} x} } } = \\
			&= \frac{a_0}{2} + \sum_{n=1}^\infty \lreck{ e^{i \frac{2\pi \cdot n}{T} x} \lr{ \frac{a_n}{2} + \frac{b_n}{2i} } + e^{-i \frac{2\pi \cdot n}{T} x} \lr{ \frac{a_n}{2} - \frac{b_n}{2i} } } = \\
			&= \frac{a_0}{2} + \sum_{n=1}^\infty e^{i \frac{2\pi \cdot n}{T} x} \lr{ \frac{a_n}{2} + \frac{b_n}{2i} } + \sum_{n=1}^\infty e^{-i \frac{2\pi \cdot n}{T} x} \lr{ \frac{a_n}{2} - \frac{b_n}{2i} } = \\
			&= \frac{1}{2} \cdot \frac{2}{T} \int_T f(x) \cdot e^{-i \frac{2\pi \cdot q}{T} x} + \sum_{n=1}^\infty e^{i \frac{2\pi \cdot n}{T} x} \lr{ \frac{a_n}{2} + \frac{b_n}{2i} } + \sum_{n=-\infty}^{-1} e^{i \frac{2\pi \cdot n}{T} x} \lr{ \frac{a_n}{2} - \frac{b_n}{2i} }
		\end{align*}
	
	Die Definition des Entwicklungskoeffizients
		\begin{align*}
			c_q &:= \frac{1}{T} \int_T f(x) \cdot e^{-i \frac{2\pi \cdot q}{T} x}
		\end{align*}
	
	ermöglicht die Zusammenfassung der drei Terme zu
		\begin{align*}
			f(x) &= \sum_{n=-\infty}^{+\infty} c_n \cdot e^{-i \frac{2\pi \cdot q}{T} x} \qquad \blacksquare
		\end{align*}


\hrule
}


\newpage


\subsection{Fourier-Transformation}
Die Fourier-Transformation ist die Darstellung einer Funktion $f(x)$ in der \underline{kontinuierlichen} Basis $\{ e^{i \frac{2\pi \cdot q}{T} x} \ | \ q \in \mathbb{R} \}$:
	\begin{alignat}{2}
		\tilde{f}(k) &= \frac{1}{\sqrt{2\pi}} \int_{-\infty}^{+\infty} \diff x \, f(x) \cdot e^{-i k x} &&\equiv \mathcal{F}\{ f(x) \}(k) \\
		f(x) &= \frac{1}{\sqrt{2\pi}} \int_{-\infty}^{+\infty} \diff k \, \tilde{f}(k) \cdot e^{i k x} &&\equiv \mathcal{F}^{-1}\{ \tilde{f}(k) \}(x)
	\end{alignat}

Wichtige Fourier-Transformationen sind:
	\begin{itemize}
		\item \textbf{Gauß-Funktion}\\
		Die Fourier-Transformation der Gauß-Funktion $f(x) = \frac{1}{\sqrt{2 \pi \sigma^2}} \cdot e^{-\frac{1}{2} \frac{(x - \mu)^2}{\sigma^2}}$ ist
			\begin{equation}
				\tilde{f}(k) = \frac{1}{\sqrt{2 \pi}} \cdot e^{- \frac{\sigma^2}{2} k^2 - i \mu k}.
			\end{equation}


		\item \textbf{Delta-Funktion}\\
		Die Fourier-Tranformation der Delta-Funktion $\delta(x-x_0)$ ist
			\begin{equation}
				\tilde{f}(k) = \frac{1}{\sqrt{2\pi}} \cdot e^{ikx_0}
			\end{equation}
		
		
		\item \textbf{Konstante}\\
		Die Fourier-Transformation der Konstanten $f(x) = A$ ist
			\begin{equation}
				\tilde{f}(k) = \sqrt{2\pi} \cdot A \cdot \delta(k).
			\end{equation}
		
		
		\item \textbf{Box-Funktion}\\
		Die Fourier-Transformation der Box-Funktion
		$f(x) = 
		\begin{cases}
			A, & -a \leq x \leq a \\
			0, & \text{sonst}
		\end{cases}$
		ist
			\begin{equation}
				\tilde{f}(k) = \frac{1}{\sqrt{2\pi a^2}} \cdot \text{sinc} \lr{ \frac{\omega}{2\pi a} }
			\end{equation}
	\end{itemize}


\newpage


{\color{violet}
\hrule
	\fbox{Herleitung}\\
	Strategie: Anwendung der Euler-McLaurin-Formel auf die komplexe Darstellung der Fourier-Reihe
		\begin{align*}
			f(x) &= \sum_{q=-\infty}^{+\infty} c_q \cdot e^{i\frac{2\pi \cdot q}{T} \cdot x} \longrightarrow \int_{-\infty}^{+\infty} \diff q \, c_q \cdot e^{i \frac{2\pi \cdot q}{T} x} = \\
			&= \int_{-\infty}^{+\infty} \diff q \lr{ \frac{1}{T} \int_T f(x) \cdot e^{-i \frac{2\pi \cdot q}{T} x} } \cdot e^{i \frac{2\pi \cdot q}{T} x} = \\
			&\stackrel{\comment{k = \frac{2\pi}{T} q}}{=} \frac{\cancel{T}}{2\pi} \int_{-\infty}^{+\infty} \diff k \lr{ \frac{1}{\cancel{T}} \int_T f(x) \cdot e^{-i \frac{2\pi \cdot q}{T} x} } \cdot e^{i \frac{2\pi \cdot q}{T} x} = \\
			&= \frac{1}{\sqrt{2\pi}} \int_{-\infty}^{+\infty} \diff k \lr{ \frac{1}{\sqrt{2\pi}} \int_T f(x) \cdot e^{-i \frac{2\pi \cdot q}{T} x} } \cdot e^{i \frac{2\pi \cdot q}{T} x} \equiv \\
			&\equiv \frac{1}{\sqrt{2\pi}} \int_{-\infty}^{+\infty} \diff k \, \tilde{f}(k) \cdot e^{i \frac{2\pi \cdot q}{T} x}
		\end{align*}

	Somit gilt
		\begin{alignat}{2}
		\tilde{f}(k) &= \frac{1}{\sqrt{2\pi}} \int_{-\infty}^{+\infty} \diff x \, f(x) \cdot e^{-i k x} &&\equiv \mathcal{F}\{ f(x) \}(k) \\
		f(x) &= \frac{1}{\sqrt{2\pi}} \int_{-\infty}^{+\infty} \diff k \, \tilde{f}(k) \cdot e^{i k x} &&\equiv \mathcal{F}^{-1}\{ \tilde{f}(k) \}(x)
	\end{alignat}

	wobei für beide Transformationen noch die abkürzenden Bezeichnungen $\mathcal{F}$ und $\mathcal{F}^{-1}$ eingeführt wurden. $\qquad \blacksquare$\\
	\\
	\fbox{Herleitung: Fourier-Transformation Gauß-Funktion}\\
	geg: $f(x) = \frac{1}{\sqrt{2 \pi \sigma^2}} \cdot e^{-\frac{1}{2} \frac{(x - \mu)^2}{\sigma^2}}$ zz: $\tilde{f}(k) = \frac{1}{\sqrt{2 \pi}} \cdot e^{- \frac{\sigma^2}{2} k^2 - i \mu k}$
		\begin{align*}
			\tilde{f}(k) &= \frac{1}{\sqrt{2\pi}} \int_{-\infty}^{+\infty} \diff x \, \frac{1}{\sqrt{2 \pi \sigma^2}} \cdot e^{-\frac{1}{2}\frac{(x-\mu)^2}{\sigma^2}} \cdot e^{-ikx} = 
			\intertext{Substituiere $u = x-\mu \Rightarrow \diff x = \diff u$:}
			&= \frac{1}{\sqrt{2\pi}} \frac{1}{\sqrt{2\pi\sigma^2}} \int_{-\infty}^{+\infty} \diff u \, e^{-\frac{1}{2}\frac{u^2}{\sigma^2}} \cdot e^{-ik(u+\mu)} = \\
			&= \frac{1}{\sqrt{2\pi}} \frac{1}{\sqrt{2\pi\sigma^2}} e^{-ik\mu} \int_{-\infty}^{+\infty} \diff u \, e^{-\frac{1}{2}\frac{1}{\sigma^2} \cdot u^2 - ik \cdot u}
			\intertext{Führe nun eine quadratische Ergänzung für den Exponenten durch, um das gängige Gauß-Integral erkennen zu können:}
				-\frac{1}{2\sigma^2} u^2 - ik u &= -\frac{1}{\sigma^2} \lr{ u^2 - i \frac{2k}{\sigma^2} u } = -\frac{1}{\sigma^2} \lr{ u^2 - i \frac{2k}{\sigma^2} u + (ik\sigma^2)^2 - (ik\sigma^2)^2 } = \\
				&= -\frac{1}{\sigma^2} \lr{ (u + ik\sigma^2)^2 - (ik\sigma^2)^2 } = -\frac{1}{\sigma^2} (u + ik\sigma^2)^2 - \frac{k^2\sigma^2}{2}
			\intertext{Einsetzen in das Integral ergibt:}
			&= \frac{1}{\sqrt{2\pi}} \frac{1}{\sqrt{2\pi\sigma^2}} e^{-ik\mu} \int_{-\infty}^{+\infty} \diff u \, e^{-\frac{1}{\sigma^2} (u + ik\sigma^2)^2 - \frac{k^2\sigma^2}{2}} = \\
			&= \frac{1}{\sqrt{2\pi}} \frac{1}{\sqrt{2\pi\sigma^2}} e^{-ik\mu} e^{-\frac{k^2\sigma^2}{2}} \int_{-\infty}^{+\infty} \diff u \, e^{-\frac{1}{2\sigma^2} (u + ik\sigma^2)^2} =
			\intertext{Substituiere $v = u + ik\sigma^2 \Rightarrow \diff u = \diff v$:}
			&= \frac{1}{\sqrt{2\pi}} \frac{1}{\sqrt{2\pi\sigma^2}} e^{-ik\mu} e^{-\frac{k^2\sigma^2}{2}} \int_{-\infty}^{+\infty} \diff v \, e^{-\frac{1}{2\sigma^2} v^2} =
			\intertext{Erkenne das Gauß-Integral $\int_{-\infty}^{+\infty} \diff x \, e^{-c\cdot x^2} = \sqrt{\pi/c}$:}
			&= \frac{1}{\sqrt{2\pi}} \frac{1}{\cancel{\sqrt{2\pi\sigma^2}}} e^{-ik\mu} e^{-\frac{k^2\sigma^2}{2}} \cdot \cancel{\sqrt{2\pi\sigma^2}} = \\
			&= \frac{1}{\sqrt{2 \pi}} \cdot e^{- \frac{\sigma^2}{2} k^2 - i \mu k} \qquad \blacksquare
		\end{align*}



\hrule
}

\newpage


\section{Stichwortverzeichnis}
\begin{xltabular}{\linewidth}{p{6cm}|X}
	\hline\hline
		Begriff & Erklärung \\
	\hline
		
	\hline\hline
\end{xltabular}
\end{document}